\documentclass[11pt,a4paper]{article}
\usepackage[T1]{fontenc}
\usepackage[utf8]{inputenc}
\usepackage{lmodern}
\usepackage{amsmath,amssymb,amsthm}
\usepackage{booktabs}
\usepackage{xcolor}
\usepackage[margin=27mm]{geometry}
\usepackage[colorlinks=true,urlcolor=blue,citecolor=blue,linkcolor=blue]{hyperref}
\newtheorem{theorem}{Theorem}
\newtheorem{proposition}[theorem]{Proposition}
\newtheorem{lemma}[theorem]{Lemma}
\newtheorem{corollary}[theorem]{Corollary}
\theoremstyle{definition}
\newtheorem{definition}[theorem]{Definition}
\newtheorem{remark}[theorem]{Remark}
\theoremstyle{plain}
\DeclareMathOperator{\Tr}{Tr}
\newcommand{\ket}[1]{\lvert#1\rangle}
\newcommand{\bra}[1]{\langle#1\rvert}
\newcommand{\norm}[1]{\lVert#1\rVert}
\newcommand{\one}{\mathbf{1}}
\newcommand{\C}{\mathbb{C}}
\setlength{\emergencystretch}{2em}
\setlength{\parindent}{0pt}
\setlength{\parskip}{0.3em plus 0.1em minus 0.1em}
\title{Exact PPT entanglement distillation from the Zhu--Wang qutrit support:\\
the one-copy optimum and the first four output dimensions}
\author{Fable \& Astra\\[0.35em]
\normalsize Research led by Jean-Philippe Bruneton\\[-0.1em]
\small Universit\'e Paris Cit\'e, France}
\date{19 November 2026}

\begin{document}
\maketitle
\begin{abstract}
Entanglement distillation converts many copies of a noisy entangled state into
fewer, perfectly entangled ones. When the conversion must be exact and the
allowed operations are those that preserve positivity of the partial transpose
(PPT), the best output from a fixed number of copies is the solution of a
semidefinite program that depends only on the support of the state. Zhu and
Wang recently exhibited a rank-three support $P$ on two qutrits for which the
asymptotic rate of exact PPT distillation lies strictly below the min-Rains
relative entropy, its standard upper bound; they bound the rate from above by
a witness but do not evaluate the one-copy program. We solve that program in
closed form. The optimal value is
$W_0(P)=\sqrt{177}/(7+\sqrt{177})=(177-7\sqrt{177})/128\approx0.65524$, so a
single copy cannot yield even one exact Bell pair. The proof reduces the
problem by a phase symmetry of $P$ to six real parameters and solves the
reduced problem exactly. Combined with the witness of Zhu and Wang, the result
fixes the largest exact maximally entangled output obtainable from $n$
copies of any state supported on $P$: its dimension is $1$, $2$, $3$ and $5$
for $n=1,2,3,4$. Four copies therefore give a five-dimensional maximally
entangled state, more than the dimension four obtained from two separate
two-copy runs. The asymptotic rate lies in
$[\log_2(1+7/\sqrt{177}),\log_2(391/250))\approx[0.6099,0.6452)$.
\end{abstract}

\section{Introduction}
Distillation is the basic test of whether the entanglement in a mixed state
is usable: Alice and Bob receive $n$ copies of a bipartite state $\rho$ and
try to produce a maximally entangled state
$\Phi_M=\frac1M\sum_{i,j=0}^{M-1}\ket{ii}\bra{jj}$ of the largest possible
Schmidt rank $M$ \cite{Bennett1996}. In the exact (zero-error) version of the
task, the output must equal $\Phi_M$ with certainty. Operations of local
quantum processing with classical communication are hard to characterize, and
a standard relaxation allows every operation that maps states with positive
partial transpose to states with positive partial transpose (PPT operations)
\cite{Rains2001}. PPT states cannot be distilled at all
\cite{Horodecki1998}, so the relaxation keeps the essential obstruction while
turning the optimization into a semidefinite program.

For exact PPT distillation from a state with support projection $P$, the
program reads
\begin{equation}\label{eq:W0}
W_0(P)=\min\bigl\{\norm{E^\Gamma}_\infty:\ P\le E\le\one\bigr\},
\end{equation}
where $E^\Gamma$ is the partial transpose and $\norm{\cdot}_\infty$ the
operator norm; $\Phi_M$ is reachable from one copy exactly when
$M\,W_0(P)\le1$, so the best integer output is
$\lfloor W_0(P)^{-1}\rfloor$ \cite{WangDuan2016,ZhuWang2026}. The asymptotic
rate $E^\infty_{0,\mathrm{PPT}}(\rho)$ is the limit of
$-\frac1n\log_2W_0(P^{\otimes n})$. It is bounded above by the min-Rains
relative entropy $R_{\min}$ \cite{WangDuan2017}, and Zhu and Wang have shown
that this bound is not tight \cite{ZhuWang2026}. Their argument uses the fact
that every feasible $E$ acts as the identity on the support, and a
range-supported witness $T$ that gives $W_0(P^{\otimes n})\ge
\norm{T^\Gamma}_1^{-n}$ for every $n$. For their explicit qutrit support the
witness bound lies strictly below $R_{\min}$.

The witness controls the problem from one side only. This note answers the
complementary question for the same support: what is $W_0(P)$ itself, and
what can actually be distilled from a few copies? The answer is exact. The
support $P$ is invariant under the phase rotations $U\otimes U$ with $U$
diagonal; averaging over them makes the optimal effect block diagonal, the
remaining six parameters decouple into a scalar inequality, and
\[
W_0(P)=\lambda_*:=\frac{\sqrt{177}}{7+\sqrt{177}}=\frac{177-7\sqrt{177}}{128}
\approx0.6552426.
\]
The value is attained by an explicit effect $E_*$ with
$\norm{E_*^\Gamma}_\infty=\lambda_*$. Because $E_*^{\otimes n}$ is feasible for
$n$ copies, the maximal output dimension $M_n$ satisfies
$M_n\ge\lfloor\lambda_*^{-n}\rfloor$, and by the witness of Zhu and Wang
$M_n<(391/250)^n$. These two bounds have the same integer part for
$n\le4$, which gives $M_1,\dots,M_4=1,2,3,5$ exactly for every state with
support $P$.

\section{Related work}\label{sec:related}
Bennett, DiVincenzo, Smolin and Wootters introduced entanglement distillation
of mixed states and its connection to quantum error correction
\cite{Bennett1996}. The Horodecki family showed that states with positive
partial transpose are undistillable, which makes the partial transpose the
natural tool for upper bounds \cite{Horodecki1998}. Rains formulated the
semidefinite relaxation by PPT operations and the corresponding upper bound on
distillable entanglement \cite{Rains2001}.

Wang and Duan gave the semidefinite programs for exact PPT distillation that
underlie the formulation \eqref{eq:W0} \cite{WangDuan2016}, and introduced the
bound later called the min-Rains relative entropy, showing that the Rains bound
is not additive \cite{WangDuan2017}. Fang, Wang, Tomamichel and Duan developed
the non-asymptotic theory of distillation under PPT operations, including the
one-shot exact setting \cite{Fang2019}. Our reformulation of the feasibility
condition $M W_0(P)\le1$ and of the integer output
$\lfloor W_0(P)^{-1}\rfloor$ follows the presentation of Zhu and Wang, who
attribute it to \cite{WangDuan2016}.

The present note depends directly on Zhu and Wang \cite{ZhuWang2026}. They
proved that every feasible effect satisfies $EP=PE=P$, constructed the
range-supported witness and the bound
$W_0(P^{\otimes n})\ge\norm{T^\Gamma}_1^{-n}$ (their Theorem~2.2), and gave the
explicit qutrit support $P$ used here together with the rational certificate
$\norm{T^\Gamma}_1<391/250$ and the lower bound
$R_{\min}\ge-\log_2(6393/10000)\approx0.6454$ (their equation~(3.9)). They do
not compute $W_0(P)$, give no achievability bound on the rate, and do not use a
symmetrization argument. The exact value of $W_0(P)$, the explicit optimal
effect, the resulting lower bound on the rate and the values of $M_1,\dots,M_4$
were not found in this source or in the other works cited here; this is a
statement about a documented search, not a certification of priority.

\section{Setting}\label{sec:setting}
Let $\{\ket0,\ket1,\ket2\}$ be the computational basis of $\C^3$ and write
$\ket{ij}=\ket i\otimes\ket j$. The support of Zhu and Wang
\cite[Sec.~3]{ZhuWang2026} is the rank-three projector
$P=\sum_{k=0}^2\ket{u_k}\bra{u_k}$ with
\begin{equation}\label{eq:support}
\ket{u_0}=\frac{8\ket{01}+7\ket{10}}{\sqrt{113}},\qquad
\ket{u_1}=\frac{7\ket{12}+8\ket{21}}{\sqrt{113}},\qquad
\ket{u_2}=\frac{\ket{20}+i\ket{02}}{\sqrt2}.
\end{equation}
For $i<j$ let $H_{ij}=\operatorname{span}\{\ket{ij},\ket{ji}\}$; we call the
three pairs $e\in\{01,12,02\}$ edges. In the ordered bases
$(\ket{ij},\ket{ji})$ the projector $P$ is block diagonal with blocks
\begin{equation}\label{eq:blocks}
P_{01}=\frac1{113}\begin{pmatrix}64&56\\56&49\end{pmatrix},\qquad
P_{12}=\frac1{113}\begin{pmatrix}49&56\\56&64\end{pmatrix},\qquad
P_{02}=\frac12\begin{pmatrix}1&i\\-i&1\end{pmatrix},
\end{equation}
which are equations (3.6)--(3.8) of \cite{ZhuWang2026}, and $P$ vanishes on
$\operatorname{span}\{\ket{00},\ket{11},\ket{22}\}$. Each block is a rank-one
projector. We denote by $p_e$ the larger diagonal entry of $P_e$ and by $c_e$
its off-diagonal entry in position $(\ket{ij},\ket{ji})$:
\[
p_{01}=p_{12}=\tfrac{64}{113},\quad p_{02}=\tfrac12,\qquad
c_{01}=c_{12}=\tfrac{56}{113},\quad c_{02}=\tfrac i2 .
\]
The partial transpose acts on the second factor,
$(\ket{ij}\bra{kl})^\Gamma=\ket{il}\bra{kj}$. Transposing the first factor
instead gives the full transpose of $E^\Gamma$, which has the same operator
norm.

The quantity of interest is $W_0(P)$ in \eqref{eq:W0}, and for $n$ copies
$W_0(P^{\otimes n})$, with the partial transpose taken over all of Bob's
qutrits. For a state $\rho$ with support $P$, let $M_n$ be the largest $M$
such that some PPT operation maps $\rho^{\otimes n}$ to $\Phi_M$ exactly. As
recalled above, $M_n=\lfloor W_0(P^{\otimes n})^{-1}\rfloor$
\cite{WangDuan2016,ZhuWang2026}; it depends on $\rho$ only through $P$.

\begin{lemma}[Fixed action {\cite[Lemma~2.1]{ZhuWang2026}}]\label{lem:fixed}
If $P$ is a projector and $P\le E\le\one$, then $EP=PE=P$. Consequently
$E=P+QEQ$ with $Q=\one-P$ and $0\le QEQ\le Q$.
\end{lemma}
\begin{proof}
For $\ket v$ in the range of $P$, $1=\bra vP\ket v\le\bra vE\ket v\le1$, so
$\bra v(\one-E)\ket v=0$. Since $\one-E\ge0$, this forces $(\one-E)\ket v=0$.
Hence $EP=P$, and $PE=P$ by taking adjoints.
\end{proof}

\section{The one-copy optimum}\label{sec:main}
\begin{theorem}\label{thm:main}
For the support \eqref{eq:support},
\[
W_0(P)=\lambda_*=\frac{\sqrt{177}}{7+\sqrt{177}}=\frac{177-7\sqrt{177}}{128}
\approx0.6552426338.
\]
The minimum is attained by the effect $E_*$ that is zero on
$\ket{00},\ket{11},\ket{22}$ and equals $P_e+t_e(\one-P_e)$ on each $H_e$,
with
\[
t_e=\frac{\lambda_*-p_e}{1-p_e},\qquad\text{that is}\qquad
\begin{aligned}
t_{01}=t_{12}&=\frac{241}{128}-\frac{113\sqrt{177}}{896}\approx0.20495,\\
t_{02}&=2\lambda_*-1=\frac{113-7\sqrt{177}}{64}\approx0.31049 .
\end{aligned}
\]
In particular one copy of any state supported on $P$ gives no exact Bell pair
under PPT operations: $M_1=\lfloor1+7/\sqrt{177}\rfloor=1$.
\end{theorem}

The proof has three steps: a symmetry reduction, a lower bound for symmetric
effects, and the verification that $E_*$ attains it.

\paragraph{Symmetry reduction.}
For real phases $\varphi=(\varphi_0,\varphi_1,\varphi_2)$ let
$U_\varphi=\sum_k e^{i\varphi_k}\ket k\bra k$. The vectors \eqref{eq:support}
are eigenvectors of $U_\varphi\otimes U_\varphi$, with eigenvalues
$e^{i(\varphi_0+\varphi_1)}$, $e^{i(\varphi_1+\varphi_2)}$ and
$e^{i(\varphi_0+\varphi_2)}$, so $(U_\varphi\otimes U_\varphi)P(U_\varphi\otimes
U_\varphi)^\dagger=P$. If $E$ is feasible, so is
$E_\varphi=(U_\varphi\otimes U_\varphi)E(U_\varphi\otimes U_\varphi)^\dagger$.
On a matrix unit, conjugation by $U_\varphi\otimes U_\varphi$ multiplies
$\ket{ij}\bra{kl}$ by $e^{i(\varphi_i+\varphi_j-\varphi_k-\varphi_l)}$, and
conjugation by $U_\varphi\otimes\bar U_\varphi$ multiplies its partial
transpose $\ket{il}\bra{kj}$ by the same phase. Hence
$E_\varphi^\Gamma=(U_\varphi\otimes\bar U_\varphi)E^\Gamma(U_\varphi\otimes\bar
U_\varphi)^\dagger$ and $\norm{E_\varphi^\Gamma}_\infty=\norm{E^\Gamma}_\infty$.
The average $\bar E=(2\pi)^{-3}\int E_\varphi\,d\varphi$ is feasible by
convexity of the constraints, and by convexity of the norm
$\norm{\bar E^\Gamma}_\infty\le\norm{E^\Gamma}_\infty$. The averaged effect
keeps exactly the matrix elements $\bra{ij}E\ket{kl}$ with
$\{i,j\}=\{k,l\}$ as multisets, because only then is the character
$\varphi_i+\varphi_j-\varphi_k-\varphi_l$ identically zero. It is therefore
enough to minimize over effects that are diagonal on $\ket{00},\ket{11},\ket{22}$
and block diagonal on the three spaces $H_e$.

For such an $E$, Lemma~\ref{lem:fixed} applied block by block gives
$E|_{H_e}=P_e+t_e(\one-P_e)$ with $t_e\in[0,1]$, because $\one-P_e$ has rank
one on $H_e$; on the remaining vectors $E\ket{ii}=d_i\ket{ii}$ with
$d_i\in[0,1]$. Conversely every choice of $t_e,d_i\in[0,1]$ is feasible. The
diagonal entries of $E|_{H_e}$ are $p_e+t_e(1-p_e)$ and
$(1-p_e)+t_ep_e$, and its off-diagonal entry is $(1-t_e)c_e$.

\paragraph{Partial transpose of a symmetric effect.}
The partial transpose fixes the diagonal matrix units $\ket{ij}\bra{ij}$ and
sends $\ket{ij}\bra{ji}$ to $\ket{ii}\bra{jj}$. Thus $E^\Gamma$ is diagonal on
the six vectors $\ket{ij}$ with $i\ne j$, with the diagonal entries of the edge
blocks, and has a $3\times3$ block $B=D+A$ on $\ket{00},\ket{11},\ket{22}$,
where $D=\operatorname{diag}(d_0,d_1,d_2)$ and $A$ is the Hermitian matrix with
zero diagonal and
\[
A_{ij}=\bra{ii}E^\Gamma\ket{jj}=\bra{ij}E\ket{ji}=(1-t_{ij})\,c_{ij}
\qquad(ij=01,12,02).
\]
The characteristic polynomial of a Hermitian $3\times3$ matrix with zero
diagonal is $x^3-s\,x-2\operatorname{Re}(A_{01}A_{12}\overline{A_{02}})$ with
$s=|A_{01}|^2+|A_{12}|^2+|A_{02}|^2$. Here $c_{01},c_{12}$ are real and
$c_{02}$ is imaginary, so the cyclic product is imaginary and the spectrum of
$A$ is $\{0,\pm\sqrt s\}$.

\paragraph{Lower bound.}
Let $\lambda=\norm{E^\Gamma}_\infty$ for a symmetric feasible $E$. If
$\lambda\ge1$ there is nothing to prove, so assume $\lambda<1$. The diagonal
entries give $\lambda\ge p_e+t_e(1-p_e)$, that is
\[
1-t_e\ge\frac{1-\lambda}{1-p_e}>0\qquad(e=01,12,02).
\]
For the unit eigenvector $\ket v$ of $A$ with eigenvalue $\sqrt s$, positivity
of $D$ gives $\lambda\ge\bra v(D+A)\ket v\ge\sqrt s$. Therefore
\[
\lambda^2\ge\sum_e|c_e|^2(1-t_e)^2\ge(1-\lambda)^2\sum_e\frac{|c_e|^2}{(1-p_e)^2}
=(1-\lambda)^2\Bigl(\tfrac{64}{49}+\tfrac{64}{49}+1\Bigr)
=(1-\lambda)^2\,\frac{177}{49},
\]
using $|c_e|/(1-p_e)=\frac{56/113}{49/113}=\frac87$ for $e=01,12$ and
$\frac{1/2}{1/2}=1$ for $e=02$. Taking square roots,
$7\lambda\ge(1-\lambda)\sqrt{177}$, i.e.\ $\lambda\ge\lambda_*$. By the
symmetry reduction the same bound holds for every feasible $E$.

\paragraph{Attainment.}
Since $\lambda_*\approx0.6552>64/113\approx0.5664$, the values $t_e$ of the
theorem lie in $(0,1)$, and $E_*$ is feasible. Its larger diagonal entry on
$H_e$ is $p_e+t_e(1-p_e)=\lambda_*$, and the smaller one,
$(1-p_e)+t_ep_e$, does not exceed it because $1-p_e\le p_e$. With
$d_i=0$ the block $B$ equals $A$, and $1-t_e=(1-\lambda_*)/(1-p_e)$ gives
$s=(1-\lambda_*)^2\cdot177/49=\lambda_*^2$. The spectrum of $E_*^\Gamma$ is
therefore contained in $[-\lambda_*,\lambda_*]$ and contains $\pm\lambda_*$,
so $\norm{E_*^\Gamma}_\infty=\lambda_*$. The closed forms follow from
$\lambda_*=\sqrt{177}(\sqrt{177}-7)/(177-49)=(177-7\sqrt{177})/128$.
Finally $\lambda_*^{-1}=1+7/\sqrt{177}\approx1.5262<2$, so $M_1=1$.
\hfill$\square$

\begin{remark}[Dual witness]\label{rem:dual}
Any density matrix $Y$ certifies a lower bound without the symmetry
reduction: by Lemma~\ref{lem:fixed}, every feasible $E$ satisfies
\[
\norm{E^\Gamma}_\infty\ge\Tr(E^\Gamma Y)=\Tr(PY^\Gamma)+\Tr(QEQ\,Y^\Gamma)
\ge\Tr(PY^\Gamma)-\Tr\bigl[(QY^\Gamma Q)_-\bigr],
\]
where $X_-$ is the negative part. Let $\ket\phi$ be the eigenvector of
$E_*^\Gamma$ for the eigenvalue $\lambda_*$ (it lies in the span of
$\ket{00},\ket{11},\ket{22}$) and
\[
Y=(1-\lambda_*)\ket\phi\bra\phi+\lambda_*\Bigl(\tfrac{64}{177}\ket{01}\bra{01}
+\tfrac{64}{177}\ket{21}\bra{21}+\tfrac{49}{177}\ket{02}\bra{02}\Bigr).
\]
The weights mirror the proof: $64/49$, $64/49$ and $1$ are the squared ratios
$|c_e|^2/(1-p_e)^2$. Numerically, the right-hand side equals $\lambda_*$ to
within $10^{-15}$, an independent confirmation of the lower bound. We have not
written out an exact evaluation of this certificate.
\end{remark}

\section{Several copies}\label{sec:copies}
\begin{corollary}\label{cor:copies}
For every state $\rho$ with support $P$ and every $n\ge1$,
\[
\Bigl(1+\frac7{\sqrt{177}}\Bigr)^{n}\le W_0(P^{\otimes n})^{-1}
<\Bigl(\frac{391}{250}\Bigr)^{n}.
\]
Consequently $M_1=1$, $M_2=2$, $M_3=3$ and $M_4=5$, while $M_5\in\{8,9\}$.
The asymptotic rate satisfies
\[
\log_2\Bigl(1+\frac7{\sqrt{177}}\Bigr)\le E^\infty_{0,\mathrm{PPT}}(\rho)
<\log_2\frac{391}{250},
\qquad\text{numerically}\quad0.60990\le E^\infty_{0,\mathrm{PPT}}(\rho)<0.64525 .
\]
\end{corollary}
\begin{proof}
The upper bound is the witness bound of Zhu and Wang,
$W_0(P^{\otimes n})\ge\norm{T^\Gamma}_1^{-n}$ with
$\norm{T^\Gamma}_1<391/250$ \cite[Theorem~2.2 and Eq.~(3.9)]{ZhuWang2026}.
For the lower bound, write $E_*=P+R$ with $R=QE_*Q\ge0$ and $PR=0$
(Lemma~\ref{lem:fixed}). Expanding $E_*^{\otimes n}$ gives $P^{\otimes n}$ plus
a sum of tensor products of copies of $P$ and $R$, each positive
semidefinite, so $P^{\otimes n}\le E_*^{\otimes n}\le\one$. Since
$(E_*^{\otimes n})^\Gamma=(E_*^\Gamma)^{\otimes n}$ and the operator norm is
multiplicative on tensor products,
$W_0(P^{\otimes n})\le\lambda_*^n$, and $\lambda_*^{-1}=1+7/\sqrt{177}$.

For the integers, $\sqrt{177}<13.305$ gives
$1+7/\sqrt{177}>1+7/13.305>1.52612$, and the two bounds become
\begin{center}
\begin{tabular}{@{}cccc@{}}
\toprule
$n$ & lower bound $\lambda_*^{-n}$ & upper bound $(391/250)^n$ & $M_n$\\
\midrule
1 & 1.5262 & 1.5640 & 1\\
2 & 2.3291 & 2.4461 & 2\\
3 & 3.5546 & 3.8257 & 3\\
4 & 5.4249 & 5.9834 & 5\\
5 & 8.2792 & 9.3580 & 8 or 9\\
\bottomrule
\end{tabular}
\end{center}
Since $M_n$ is an integer with $\lambda_*^{-n}\le W_0(P^{\otimes n})^{-1}$,
we have $M_n\ge\lfloor\lambda_*^{-n}\rfloor$, and $M_n<(391/250)^n$. For
$n\le4$ the lower bound is at least $M$ and the upper bound is at most $M+1$
for the value $M$ listed in the table, which fixes $M_n=M$. The checker
verifies these inequalities in rational arithmetic. For the rate,
$\frac1n\log_2M_n\ge\frac1n\log_2\lfloor\lambda_*^{-n}\rfloor\to
\log_2\lambda_*^{-1}$, and the upper bound is that of \cite{ZhuWang2026}.
\end{proof}

Two features of these numbers deserve comment. First, one copy yields
nothing while two copies yield an exact Bell pair; the gain comes entirely
from the tensor structure, since $E_*^{\otimes2}$ has norm
$\lambda_*^2\approx0.4293<1/2$. Second, the exact output is superadditive in
the number of copies: $M_3=3>M_1M_2=2$ and $M_4=5>M_2^2=4$, so four copies
processed jointly give a five-dimensional maximally entangled state, whereas
two independent two-copy runs give dimension four. These statements hold for
every state supported on $P$, because $W_0$ depends only on the support.

The lower bound $\log_2(1+7/\sqrt{177})\approx0.6099$ on the rate comes from
tensor powers of the one-copy optimum. If $W_0$ were multiplicative on this
support, it would be the exact rate; the witness value
$\log_2(391/250)\approx0.6452$ would be the exact rate if the witness were
attained. The gap between them, about $0.035$ bits per copy, is what the
present results leave open. Both values lie below the min-Rains lower bound
$-\log_2(6393/10000)\approx0.6454$ of \cite{ZhuWang2026}.

\section{Verification}\label{sec:verification}
The theorem and the corollary are proved above. The public
\href{https://github.com/jpbruneton/open_quantum_problems/blob/main/research/entries/exact-ppt-distillation-qutrit/checks/verify_ppt_distillation.py}{Python checker}
and its
\href{https://github.com/jpbruneton/open_quantum_problems/blob/main/research/entries/exact-ppt-distillation-qutrit/data/verification.json}{recorded output}
complement the proofs in four steps. With exact symbolic arithmetic (SymPy) it
builds $P$ from \eqref{eq:support}, checks the blocks \eqref{eq:blocks}, builds
$E_*$, checks feasibility, and verifies that $E_*^\Gamma$ has diagonal
entries at most $\lambda_*$ and that its $3\times3$ block has characteristic
polynomial $x(x^2-\lambda_*^2)$. In floating point it then tests the lower
bound in three ways: $20\,000$ random feasible effects $E=P+QXQ$ with
$0\le X\le\one$ (smallest norm found $\approx0.779$), a multistart
minimization over the six-parameter symmetric family (minimum $\lambda_*$ to
machine precision), and the dual witness of Remark~\ref{rem:dual}. The
first and third tests do not use the symmetry reduction. Finally it
verifies the integer bounds of Corollary~\ref{cor:copies} with rational
numbers. No input certificate is needed: the support is given by the rational
data \eqref{eq:support}.

\section{Scope and limitations}\label{sec:scope}
The exact value of $W_0(P)$ and the optimal effect are proved analytically.
The values $M_1,\dots,M_4$ and the rate interval follow from that proof
combined with the witness bound of Zhu and Wang, whose certified inequality
$\norm{T^\Gamma}_1<391/250$ we take from their paper and do not re-derive.
The numerical tests are a consistency check, not part of the proof. The
results concern this particular support; the method extends to supports with
the same phase symmetry, where the reduced problem again involves a $3\times3$
block, but the cancellation of the cyclic product used here depends on the
phase $i$ in $\ket{u_2}$. Whether $W_0(P^{\otimes n})=\lambda_*^n$ for all
$n$, the value $M_5\in\{8,9\}$, and the exact asymptotic rate remain open. The
bibliographic comparison of Section~\ref{sec:related} is a targeted reading of
the cited sources and does not establish priority.

\begingroup\small
\begin{thebibliography}{99}
\setlength{\itemsep}{2pt}
\bibitem{Bennett1996}
C.~H. Bennett, D.~P. DiVincenzo, J.~A. Smolin and W.~K. Wootters,
\emph{Mixed-state entanglement and quantum error correction},
Phys. Rev. A \textbf{54}, 3824 (1996);
\href{https://doi.org/10.1103/PhysRevA.54.3824}{doi:10.1103/PhysRevA.54.3824};
\href{https://arxiv.org/abs/quant-ph/9604024}{arXiv:quant-ph/9604024}.

\bibitem{Horodecki1998}
M. Horodecki, P. Horodecki and R. Horodecki,
\emph{Mixed-state entanglement and distillation: is there a ``bound''
entanglement in nature?},
Phys. Rev. Lett. \textbf{80}, 5239 (1998);
\href{https://doi.org/10.1103/PhysRevLett.80.5239}{doi:10.1103/PhysRevLett.80.5239};
\href{https://arxiv.org/abs/quant-ph/9801069}{arXiv:quant-ph/9801069}.

\bibitem{Rains2001}
E.~M. Rains,
\emph{A semidefinite program for distillable entanglement},
\href{https://arxiv.org/abs/quant-ph/0008047}{arXiv:quant-ph/0008047} (2000).

\bibitem{WangDuan2016}
X. Wang and R. Duan,
\emph{Improved semidefinite programming upper bound on distillable
entanglement},
Phys. Rev. A \textbf{94}, 050301 (2016);
\href{https://doi.org/10.1103/PhysRevA.94.050301}{doi:10.1103/PhysRevA.94.050301};
\href{https://arxiv.org/abs/1601.07940}{arXiv:1601.07940}.

\bibitem{WangDuan2017}
X. Wang and R. Duan,
\emph{Nonadditivity of Rains' bound for distillable entanglement},
Phys. Rev. A \textbf{95}, 062322 (2017);
\href{https://doi.org/10.1103/PhysRevA.95.062322}{doi:10.1103/PhysRevA.95.062322};
\href{https://arxiv.org/abs/1605.00348}{arXiv:1605.00348}.

\bibitem{Fang2019}
K. Fang, X. Wang, M. Tomamichel and R. Duan,
\emph{Non-asymptotic entanglement distillation},
IEEE Trans. Inf. Theory \textbf{65}, no.~10 (2019);
\href{https://doi.org/10.1109/TIT.2019.2914688}{doi:10.1109/TIT.2019.2914688};
\href{https://arxiv.org/abs/1706.06221}{arXiv:1706.06221}.

\bibitem{ZhuWang2026}
C. Zhu and X. Wang,
\emph{The min-Rains relative entropy is not tight for exact PPT entanglement
distillation},
\href{https://arxiv.org/abs/2608.12135}{arXiv:2608.12135v1} (2026).
\end{thebibliography}
\endgroup
\end{document}
